\section{De la demo a la productividad: un marco de aceptación}

En este episodio aparece una y otra vez una pregunta de evaluación implícita: ¿en qué momento un robot puede considerarse utilizable? Gao Jiyang (高继扬) describe las etapas de capacidad desde la demo en video, la demo en oficina y la demo en entornos reales (\emph{wild demo}) hasta los escenarios de productividad. Esta división es más útil que fijarse solo en las métricas del modelo, porque lo que el robot termina por enfrentar no es una tabla de clasificación, sino clientes reales y tareas reales.

\subsection{Cuatro niveles de aceptación}

Aquí se organizan las formas de demo mencionadas en la entrevista en una tabla de aceptación de cuatro niveles. El propósito de la tabla no es inventar términos nuevos, sino recordar al lector que la capacidad de un robot debe pasar de «verlo tener éxito una vez» a «generar valor de forma continua para el cliente».

\begin{longtable}{p{0.22\textwidth}p{0.34\textwidth}p{0.35\textwidth}}
\toprule
Etapa & Qué demuestra & Qué no demuestra \\
\midrule
Demo in Video & Es visualmente atractiva; demuestra que una acción puede captarse en video & No demuestra estabilidad, reproducibilidad ni adaptación al sitio. \\
Demo in Office & Puede presentarse en vivo en el entorno controlado de la empresa & No demuestra que funcione en distintos lugares, con distintos objetos ni en los procesos de distintos clientes. \\
Demo in the Wild & Puede desplegarse y presentarse con clientes, en ferias o en escenarios en otros países & No equivale a operación prolongada sin supervisión ni a un bajo costo de mantenimiento. \\
Escenario de productividad & El cliente está dispuesto a usarlo de forma continua y a pagar por el valor que recibe & Aún requiere validar la entrega a escala, el mantenimiento y el ciclo cerrado de datos. \\
\bottomrule
\end{longtable}

\begin{warningbox}{El sesgo en la percepción de los robots que generan los videos}
La idea que el público general tiene de los robots suele provenir de videos cortos. Los videos amplifican los fragmentos exitosos y minimizan el número de fallas, el tiempo de preparación, las limitaciones del escenario y el costo de mantenimiento. Por ello, para evaluar a una empresa de robótica no basta con ver las demos: hay que observar los ciclos, los envíos, la recompra de los clientes, el mantenimiento y el retorno de datos.
\end{warningbox}

\subsection{El sitio del cliente es la evaluación definitiva}

Gao Jiyang dice que, si un cliente opina que el producto no es bueno, él acude de inmediato al sitio del cliente para resolverlo y convierte cada problema individual en la solución sistemática de toda una clase de problemas. Esta práctica corresponde a la forma real de evaluar la inteligencia corporizada: no es una calificación única en una tabla de clasificación fuera de línea, sino exponer problemas una y otra vez en el sitio del cliente, resolverlos, consolidar los procesos y mejorar el producto.

\begin{importantbox}{La triple validación de los escenarios de productividad}
Para que una capacidad robótica entre en un escenario de productividad debe superar, como mínimo, una triple validación: que el cliente esté dispuesto a pagar, que el sistema complete las tareas de forma estable y que el sistema de mantenimiento y el ciclo cerrado de datos absorban los problemas de manera continua.
\end{importantbox}

\subsection{Resumen de la sección}

Para pasar de la demo a la productividad, una empresa de robótica debe superar cuatro umbrales: que el sistema sea reproducible, desplegable, mantenible y que el cliente pague por él. La verdadera capacidad del modelo tiene que redefinirse a la luz de estos umbrales.
